\section{Proofs supporting the one-dimensional reduction}
\label{app:local-reduction}

\begin{proof}[Proof of \Cref{lem:boundary-convergence}]
Put
\[
  q_n=\pc(r_n),
  \qquad
  p_0=\pc(r_0).
\]
By continuity of $\pc$ on $K$, we have $q_n\to p_0$.  Since $\pc(K)$ is a compact subset of $(0,1)$, choose a
compact interval $P\subset(0,1)$ whose interior contains $\pc(K)$.  Then
$q_n\in P$ for every $n$, and $|p_n-q_n|\to0$ implies, after discarding
finitely many terms, that $p_n\in P$ as well.
The role of $P$ is only to keep the logarithms in $I_p$ uniformly away from the
singular endpoints $p=0,1$.

Since $(\widetilde{\W}_0,\delta_\cut)$ is compact, every subsequence of
$(W_n)$ has a further cut-convergent subsequence; so it suffices to show that
every cut-convergent subsequence has limit $r_0$.
Passing to such a subsequence and relabelling it, we assume that $W_n\to W_*$ in cut distance.
We will prove that $W_*\equiv r_0$.
First note that the homomorphism densities
are continuous in the cut metric, so
\begin{equation}
  \label{eq:boundary-limit-feasibility}
  t(H,W_*)\ge r_0^m.
\end{equation}

We next compare the relative-entropy functionals with nearby values of the
parameter $p$.  For $p,q\in P$, the identity
\[
  I_p(W)
  =
  -2s(W)-\log(1-p)+e(W)\log\frac{1-p}{p}
\]
gives
\[
  I_p(W)-I_q(W)
  =
  \log\frac{1-q}{1-p}
  +e(W)\log\frac{(1-p)q}{p(1-q)}.
\]
Let
\[
  A(p')=-\log(1-p'),
  \qquad
  B(p')=\log\frac{1-p'}{p'}.
\]
Then, the preceding identity can be written as
\[
  I_p(W)-I_q(W)=A(p)-A(q)+e(W)\bigl(B(p)-B(q)\bigr).
\]
Note that as $P$ is compact, both $A$ and $B$
are uniformly continuous on $P$; the singularities at $0$ and $1$ are now a
positive distance away.
Since every graphon satisfies
$0\le e(W)\le1$, we have, for all $p,q\in P$,
\[
  |I_p(W)-I_q(W)| = |A(p)-A(q)+e(W)(B(p)-B(q))|
  \le |A(p)-A(q)| + |B(p)-B(q)|.
\]
Thus, by the uniform continuity of $A$ and $B$ on $P$,
\begin{equation}\label{eq:entropy-parameter-continuity}
\sup_{W\in\W_0}|I_p(W)-I_q(W)|\to0
\quad\text{as } |p-q|\to0,\quad p,q\in P.
\end{equation}


For fixed $p_0$, lower semicontinuity of $I_{p_0}$ in the cut metric gives
$I_{p_0}(W_*)\le\liminf_n I_{p_0}(W_n)$.  Since
\eqref{eq:entropy-parameter-continuity} and $q_n\to p_0$ give
$|I_{p_0}(W_n)-I_{q_n}(W_n)|\to0$, we obtain
\begin{equation}\label{eq:entropy-moving-parameter-lsc}
  I_{p_0}(W_*)
  \le
  \liminf_{n\to\infty} I_{q_n}(W_n).
\end{equation}

On the other hand, the constant graphon $W\equiv r_n$ is feasible for
$(p_n,r_n)$.  Since $W_n$ is an optimizer,
\[
  I_{p_n}(W_n)\le J_{p_n}(r_n).
\]
Using \eqref{eq:entropy-parameter-continuity} again, now with $p=p_n$ and $q=q_n$,
we obtain
\[
  I_{q_n}(W_n)
  \le
  J_{p_n}(r_n)+o(1).
\]
Since $p_n-q_n\to0$, $q_n\to p_0$, and $r_n\to r_0$, continuity of
$(p,u)\mapsto J_p(u)$ on $(0,1)\times[0,1]$, hence uniform continuity on
compact subsets, gives
\[
  J_{p_n}(r_n)=J_{q_n}(r_n)+o(1)\to J_{p_0}(r_0).
\]
Therefore
\begin{equation}\label{eq:boundary-cost-limsup}
  \limsup_{n\to\infty} I_{q_n}(W_n)
  \le
  J_{p_0}(r_0).
\end{equation}

Equations \eqref{eq:boundary-limit-feasibility}, \eqref{eq:entropy-moving-parameter-lsc}, and \eqref{eq:boundary-cost-limsup} show
that $W_*$ is feasible at $(p_0,r_0)$ and has $I_{p_0}$-value no larger than the
constant graphon:
\[
  I_{p_0}(W_*)\le J_{p_0}(r_0).
\]
By condition (M2) of \Cref{thm:scalar-lz-boundary},
$(r_0^d,J_{p_0}(r_0))$ lies on the convex minorant of
$x\mapsto J_{p_0}(x^{1/d})$.  Hence
Theorem~\ref{thm:lz-criterion} gives that the constant graphon $r_0$
is the unique minimizer at $(p_0,r_0)$.  Thus,  $W_*\equiv r_0$.  Since every subsequence of $(W_n)$ has a further
cut-convergent subsequence, and every such limit equals $r_0$, the whole
sequence converges to $r_0$ in the quotient cut metric.  Finally, edge
density is continuous in the cut metric, so $e(W_n)\to r_0$.
\end{proof}

\begin{proof}[Proof of \Cref{cor:scalar-reduction}]
Fix $0<\rho<\rho_0$, write $\eta=\eta(\rho)$, and take
$r\in K$, $p\in(\pc(r)-\eta,\pc(r))$, and an optimizer $W^*$ attaining
$\Phi_H(p,r)$.  By \Cref{lem:active-constraint}, $t(H,W^*)=r^m$, and by
\Cref{lem:fixed-density-bipodality},
$\eps^*=e(W^*)\in(r-\rho,r)$.  Thus, $W^*$ is admissible in the variational
problem defining $\sigma_H(\eps^*,r^m)$, and hence
\[
  s(W^*)\le\sigma_H(\eps^*,r^m).
\]
Hence, by \eqref{eq:relative-entropy-identity},
\[
  \mathcal F_{p,r}(\eps^*)\le I_p(W^*)=\Phi_H(p,r).
\]
On the other hand, $\widehat W_{\eps^*,r}$ has edge density $\eps^*$ and
$H$-density $r^m$, so it is feasible for the variational problem defining
$\Phi_H(p,r)$ in \eqref{eq:graphon-variational-problem}.  Using
\eqref{eq:reduced-objective-attainment}, we get
\[
  \Phi_H(p,r)\le I_p(\widehat W_{\eps^*,r})=\mathcal F_{p,r}(\eps^*).
\]
Thus, the preceding inequalities are all equalities, and in particular,
\[
  \mathcal F_{p,r}(\eps^*)=\Phi_H(p,r),
\]
and
\[
  s(W^*)=\sigma_H(\eps^*,r^m).
\]
To see this, suppose that the entropy inequality
$s(W^*)\le\sigma_H(\eps^*,r^m)$ were strict.  Then the inequality
$\mathcal F_{p,r}(\eps^*)\le I_p(W^*)$ would also be strict.  Therefore,
$W^*$ is an entropy maximizer under the fixed constraints
$e(W)=\eps^*$ and $t(H,W)=r^m$.  By the uniqueness part of
\Cref{thm:krrs-analytic-extension}, $W^*$ agrees with $\widehat W_{\eps^*,r}$ up to relabelling.

For any $\eps\in(r-\rho,r)$, the graphon $\widehat W_{\eps,r}$ is feasible for the
variational problem defining $\Phi_H(p,r)$ in \eqref{eq:graphon-variational-problem}, so
\[
  \mathcal F_{p,r}(\eps)=I_p(\widehat W_{\eps,r})\ge\Phi_H(p,r)
  =\mathcal F_{p,r}(\eps^*).
\]
Therefore, $\eps^*$ minimizes $\mathcal F_{p,r}$ on $(r-\rho,r)$, and
\eqref{eq:edge-density-minimization} follows.  The same inequality gives the converse
statement: if $\eps$ minimizes $\mathcal F_{p,r}$ on $(r-\rho,r)$, then
$I_p(\widehat W_{\eps,r})=\mathcal F_{p,r}(\eps)=\mathcal F_{p,r}(\eps^*)
=\Phi_H(p,r)$, and $\widehat W_{\eps,r}$ is feasible, so it is an optimizer
attaining $\Phi_H(p,r)$.
Finally, distinct $\eps$ give graphons of distinct edge density, hence
inequivalent under relabelling, which gives the asserted bijection.
\end{proof}